\documentclass[aspectratio=169,11pt]{beamer}
\usepackage[utf8]{inputenc}
\usepackage[T1]{fontenc}
\usepackage[spanish,es-nodecimaldot]{babel}
\usepackage{tikz}
\usetikzlibrary{positioning,arrows.meta,fit,calc}
\usepackage{booktabs}
\usepackage{array}

\definecolor{ink}{HTML}{0B1320}
\definecolor{panel}{HTML}{142338}
\definecolor{panel2}{HTML}{1C314A}
\definecolor{mint}{HTML}{5EEAD4}
\definecolor{orange}{HTML}{FDBA74}
\definecolor{cream}{HTML}{F8FAFC}
\definecolor{muted}{HTML}{A9B8C8}
\definecolor{red}{HTML}{FB7185}

\setbeamercolor{normal text}{fg=cream,bg=ink}
\setbeamercolor{background canvas}{bg=ink}
\setbeamercolor{frametitle}{fg=cream,bg=ink}
\setbeamercolor{title}{fg=cream}
\setbeamercolor{subtitle}{fg=mint}
\setbeamercolor{structure}{fg=mint}
\setbeamercolor{block title}{fg=ink,bg=mint}
\setbeamercolor{block body}{fg=cream,bg=panel}
\setbeamercolor{alerted text}{fg=orange}
\setbeamertemplate{navigation symbols}{}
\setbeamertemplate{blocks}[rounded][shadow=false]
\setbeamertemplate{itemize item}{\color{mint}\small$\blacktriangleright$}
\setbeamertemplate{footline}{%
  \hfill\color{muted}\scriptsize\insertframenumber\, / \,\inserttotalframenumber\hspace{0.35cm}\vspace{0.2cm}}
\setbeamertemplate{frametitle}{\vspace{0.35cm}\hspace{0.1cm}\insertframetitle\par\vspace{0.12cm}\hrule\vspace{0.18cm}}
\setlength{\parskip}{3pt}
\newcommand{\tagline}[1]{{\color{mint}\footnotesize\bfseries #1}}
\newcommand{\bigaccent}[1]{{\color{orange}\fontsize{28}{30}\selectfont\bfseries #1}}
\newcommand{\pill}[1]{\tikz[baseline=(x.base)]\node[fill=panel2,rounded corners=3pt,inner xsep=7pt,inner ysep=3pt](x){\color{mint}\scriptsize\bfseries #1};}

\title{HayPaComer}
\subtitle{La nevera que sabe qué hay, cuánto queda y qué puedes cocinar}
\author{Proyecto de patrones de diseño en Java}
\institute{\tagline{DEL INVENTARIO A UNA DECISIÓN REAL}}
\date{Septiembre de 2026}

\begin{document}

{\setbeamertemplate{footline}{}
\begin{frame}[plain]
\begin{tikzpicture}[remember picture,overlay]
\fill[panel] (current page.south west) rectangle ([xshift=4.4cm]current page.north west);
\fill[mint] ([xshift=4.4cm]current page.south west) rectangle ([xshift=4.48cm]current page.north west);
\node[anchor=north west] at ([xshift=0.7cm,yshift=-0.7cm]current page.north west) {\color{orange}\fontsize{42}{44}\selectfont\bfseries 01};
\node[anchor=south west] at ([xshift=0.7cm,yshift=0.7cm]current page.south west) {\color{muted}\scriptsize JAVA  ·  ESP32  ·  IA};
\end{tikzpicture}
\vspace{0.8cm}
\hspace{4.9cm}\begin{minipage}{0.62\textwidth}
{\tagline{PROYECTO DE SOFTWARE + HARDWARE}}\\[0.25cm]
{\usebeamerfont{title}\color{cream}\fontsize{35}{39}\selectfont\bfseries HayPaComer}\\[0.18cm]
{\color{mint}\Large La nevera que sabe qué hay, cuánto queda y qué puedes cocinar.}\\[0.55cm]
{\color{muted}\large Una aplicación Java conectada a una nevera instrumentada para reducir desperdicio y decidir mejor cada comida.}
\end{minipage}
\end{frame}}

\begin{frame}{El problema no es cocinar}
\tagline{TODOS LOS DÍAS, LA MISMA FRICCIÓN}
\vspace{0.35cm}
\begin{columns}[T,onlytextwidth]
\column{0.33\textwidth}
\begin{block}{\color{orange}¿Qué hay?}
La comida está repartida, escondida y nadie tiene una respuesta confiable.
\end{block}
\column{0.33\textwidth}
\begin{block}{\color{orange}¿Qué se vence?}
Las sobras y productos abiertos desaparecen de la memoria de la casa.
\end{block}
\column{0.33\textwidth}
\begin{block}{\color{orange}¿Alcanza?}
Una receta dice “200 g”, pero en casa solo quedan 80 g.
\end{block}
\end{columns}
\vfill
\begin{center}
{\color{cream}\Large\bfseries Se compra de más. Se desperdicia. Se improvisa.}\\[0.18cm]
{\color{muted}\large La tecnología suele registrar productos; no ayuda a tomar la decisión de esta noche.}
\end{center}
\end{frame}

\begin{frame}{La idea en una frase}
\begin{center}
\vspace{0.4cm}
{\color{mint}\fontsize{24}{28}\selectfont\bfseries HayPaComer convierte la nevera en una decisión.}
\vspace{0.45cm}
\begin{tikzpicture}[node distance=0.55cm,>=Latex]
\node[draw=mint,very thick,rounded corners=8pt,fill=panel,minimum width=2.7cm,minimum height=1cm] (a) {\color{cream}\bfseries Datos reales};
\node[draw=orange,very thick,rounded corners=8pt,fill=panel,minimum width=2.7cm,minimum height=1cm,right=of a] (b) {\color{cream}\bfseries Java decide};
\node[draw=mint,very thick,rounded corners=8pt,fill=panel,minimum width=2.7cm,minimum height=1cm,right=of b] (c) {\color{cream}\bfseries Cocina mejor};
\draw[->,very thick,mint] (a)--(b);
\draw[->,very thick,orange] (b)--(c);
\end{tikzpicture}
\vspace{0.55cm}
\begin{columns}[T,onlytextwidth]
\column{0.5\textwidth}
\pill{INVENTARIO VIVO}\\[-1pt]
Qué hay, cuánto pesa, de quién es y cuándo vence.
\column{0.5\textwidth}
\pill{COCINA AHORA}\\[-1pt]
Una sugerencia posible en 5, 15 o 30 minutos.
\end{columns}
\end{center}
\end{frame}

\begin{frame}{¿Cómo se usa en la vida real?}
\tagline{UN FLUJO DE 30 SEGUNDOS}
\vspace{0.35cm}
\begin{enumerate}
\item Abrir la app: aparecen sobras, vencimientos y alertas.
\item Elegir: ``15 minutos, una persona, Juan''.
\item Pesar el ingrediente que se va a utilizar.
\item Confirmar: alcanza, se reduce la porción o se sustituye.
\item Cocinar: el sistema descuenta lo usado y actualiza el mercado.
\end{enumerate}
\vfill
\begin{center}
\begin{tikzpicture}[node distance=0.22cm]
\node[fill=panel2,rounded corners=5pt,inner sep=8pt] (x) {\shortstack{\color{cream}\bfseries 18:40\\[-2pt]\color{mint}\scriptsize Abrir nevera}};
\node[fill=panel2,rounded corners=5pt,inner sep=8pt,right=of x] (y) {\shortstack{\color{cream}\bfseries 18:41\\[-2pt]\color{mint}\scriptsize Pesar pollo}};
\node[fill=panel2,rounded corners=5pt,inner sep=8pt,right=of y] (z) {\shortstack{\color{cream}\bfseries 18:42\\[-2pt]\color{mint}\scriptsize Sustituir}};
\node[fill=orange,rounded corners=5pt,inner sep=8pt,right=of z] (w) {\shortstack{\color{ink}\bfseries 18:55\\[-2pt]\scriptsize Comer}};
\draw[->,thick,mint] (x)--(y);\draw[->,thick,mint] (y)--(z);\draw[->,thick,orange] (z)--(w);
\end{tikzpicture}
\end{center}
\end{frame}

\begin{frame}{La pesa: el salto de “hay” a “alcanza”}
\tagline{UNA CELDA DE CARGA · DOS MODOS}
\begin{columns}[T,onlytextwidth]
\column{0.47\textwidth}
\begin{block}{Modo NEVERA}
La bandeja mide el stock de un producto asignado. Si pasa de 842 g a 650 g, Java registra un consumo real de 192 g.
\end{block}
\begin{block}{Modo COCINA}
Se tara el recipiente, se pone el ingrediente y se compara el peso con la receta.
\end{block}
\column{0.47\textwidth}
\begin{center}
\begin{tikzpicture}[scale=0.9]
\draw[fill=panel2,draw=mint,very thick,rounded corners=5pt] (0,0) rectangle (4.5,2.5);
\draw[fill=orange!80!ink,draw=orange,thick] (0.8,0.65) rectangle (3.7,1.15);
\draw[fill=cream!80!ink,draw=cream] (1.3,1.15) rectangle (3.2,1.9);
\node at (2.25,0.9) {\color{ink}\bfseries HX711};
\node at (2.25,1.52) {\color{ink}\scriptsize ingrediente};
\draw[->,very thick,mint] (2.25,2.5)--(2.25,3.1);
\node[fill=mint,rounded corners=3pt,inner sep=4pt] at (2.25,3.35) {\color{ink}\bfseries 132 g};
\end{tikzpicture}
\vspace{0.25cm}
\color{muted}\small ESP32 + HX711 + celda de carga\\
\color{muted}\small tara + lectura estable + gramos reales
\end{center}
\end{columns}
\end{frame}

\begin{frame}{Cuando no alcanza, no se rompe la receta}
\tagline{MOTOR DE CANTIDADES Y SUSTITUCIONES}
\begin{center}
\begin{tikzpicture}[node distance=0.35cm,>=Latex]
\node[fill=panel2,rounded corners=6pt,minimum width=3cm,minimum height=0.8cm] (r) {\color{cream}\bfseries Receta pide 200 g};
\node[fill=panel2,rounded corners=6pt,minimum width=3cm,minimum height=0.8cm,right=1cm of r] (m) {\color{orange}\bfseries Se pesan 80 g};
\node[fill=mint,rounded corners=6pt,minimum width=3cm,minimum height=0.8cm,right=1cm of m] (q) {\color{ink}\bfseries Java evalúa};
\draw[->,very thick,mint] (r)--(m);\draw[->,very thick,mint] (m)--(q);
\end{tikzpicture}
\vspace{0.45cm}
\end{center}
\begin{columns}[T,onlytextwidth]
\column{0.32\textwidth}
\begin{block}{1 · Reducir}
La receta baja de 2 porciones a 1.
\end{block}
\column{0.32\textwidth}
\begin{block}{2 · Sustituir}
Se pesa atún, huevo u otra alternativa permitida.
\end{block}
\column{0.32\textwidth}
\begin{block}{3 · Faltante}
La cantidad exacta cae a la lista de mercado.
\end{block}
\end{columns}
\vfill
\begin{center}\color{orange}\bfseries La IA sugiere. Java valida los gramos.\end{center}
\end{frame}

\begin{frame}{Hardware que sí se puede construir}
\tagline{EL ESP32 MIDE; JAVA INTERPRETA}
\begin{columns}[T,onlytextwidth]
\column{0.52\textwidth}
\begin{block}{Componentes}
\begin{itemize}
\item \textcolor{mint}{ESP32}: WiFi y lectura.
\item \textcolor{mint}{Reed + imán}: puerta.
\item \textcolor{mint}{DS18B20}: temperatura.
\item \textcolor{mint}{HX711 + celda}: gramos reales.
\item \textcolor{mint}{LED + buzzer}: alerta local.
\end{itemize}
\end{block}
\column{0.45\textwidth}
\begin{block}{Flujo}
\begin{center}
\color{cream}\bfseries Sensores\\[0.12cm]
\color{mint}\Large $\downarrow$\\[0.12cm]
\color{cream}\bfseries ESP32\\[0.12cm]
\color{mint}\Large $\downarrow$\\[0.12cm]
\color{orange}\bfseries Java\\[0.12cm]
\color{mint}\Large $\downarrow$\\[0.12cm]
\color{cream}\bfseries App
\end{center}
\end{block}
\end{columns}
\end{frame}

\begin{frame}{La IA sin humo}
\tagline{IA COMO AYUDANTE, NO COMO AUTORIDAD}
\begin{columns}[T,onlytextwidth]
\column{0.32\textwidth}
\pill{SUGERIR}\vspace{0.1cm}\\
Genera platos usando inventario, tiempo, preferencias y gramos medidos.
\column{0.32\textwidth}
\pill{INTERPRETAR}\vspace{0.1cm}\\
Lee una etiqueta o recibo y propone producto, fecha y cantidad.
\column{0.32\textwidth}
\pill{EXPLICAR}\vspace{0.1cm}\\
Redacta por qué conviene consumir algo primero o sustituirlo.
\end{columns}
\vfill
\begin{block}{Regla de oro}
La IA puede decir ``podrías usar atún''. Solo Java, con reglas y mediciones, decide si la cantidad alcanza. El usuario confirma antes de descontar.
\end{block}
\vspace{0.25cm}
\begin{center}\color{muted}\small Si la API falla, el motor de reglas mantiene la aplicación funcionando.\end{center}
\end{frame}

\begin{frame}{Patrones que aparecen naturalmente}
\tagline{NO USAMOS PATRONES PARA DECORAR: RESUELVEN CAMBIOS REALES}
\begin{columns}[T,onlytextwidth]
\column{0.48\textwidth}
\begin{itemize}
\item \textcolor{mint}{Facade}: una entrada para consultar y cocinar.
\item \textcolor{mint}{Adapter}: ESP32, OCR e IA hablan con Java.
\item \textcolor{mint}{Bridge}: sensor físico separado de interpretación.
\item \textcolor{mint}{Proxy}: hardware offline y alimentos privados.
\item \textcolor{mint}{Composite}: nevera → zona → bandeja → alimento.
\end{itemize}
\column{0.48\textwidth}
\begin{itemize}
\item \textcolor{orange}{Builder}: construir plato con restricciones.
\item \textcolor{orange}{Strategy}: gramos, unidades o sustituciones.
\item \textcolor{orange}{Decorator}: vencido, sobra, en riesgo.
\item \textcolor{orange}{Factory Method}: eventos y tipos de alimento.
\item \textcolor{orange}{Flyweight}: metadatos compartidos.
\end{itemize}
\end{columns}
\end{frame}

\begin{frame}{Demo: de la nevera al plato}
\tagline{UNA HISTORIA QUE SE PUEDE MOSTRAR}
\begin{enumerate}
\item La puerta permanece abierta 40 segundos → LED rojo y buzzer.
\item La leche baja de 842 g a 650 g → stock actualizado en 192 g.
\item Se elige ``arroz con pollo, 15 minutos, 2 personas''.
\item La receta pide 200 g de pollo; la pesa confirma solo 80 g.
\item HayPaComer propone reducir a 1 porción o sustituir con atún.
\item Se pesan 130 g de atún → Java valida, descuenta y actualiza mercado.
\item La app recomienda sopa de ayer para la siguiente comida.
\end{enumerate}
\vfill
\begin{center}{\color{mint}\Large\bfseries El sistema no adivina la nevera: la mide.}\end{center}
\end{frame}

\begin{frame}{Qué hace diferente a HayPaComer}
\tagline{NO ES UNA APP DE RECETAS}
\begin{columns}[T,onlytextwidth]
\column{0.32\textwidth}
\begin{block}{Recetas comunes}
Parten de ingredientes ideales y cantidades perfectas.
\end{block}
\column{0.32\textwidth}
\begin{block}{Nevera inteligente cara}
Detecta alimentos, pero exige un electrodoméstico costoso.
\end{block}
\column{0.32\textwidth}
\begin{block}{HayPaComer}
Mide cantidades reales, acepta faltantes y funciona con hardware accesible.
\end{block}
\end{columns}
\vfill
\begin{center}
{\color{orange}\fontsize{22}{25}\selectfont\bfseries Menos desperdicio.\\ Más decisiones correctas.}
\end{center}
\end{frame}

{\setbeamertemplate{footline}{}
\begin{frame}[plain]
\begin{center}
\vspace{1cm}
{\color{mint}\fontsize{34}{38}\selectfont\bfseries HayPaComer}\\[0.35cm]
{\color{cream}\Large La pregunta ya no es ``¿qué habrá?''}\\
{\color{orange}\Large\bfseries Es: ``¿qué puedo hacer con lo que tengo?''}\\[1cm]
\pill{JAVA}\hspace{0.2cm}\pill{ESP32}\hspace{0.2cm}\pill{HX711}\hspace{0.2cm}\pill{IA}\hspace{0.2cm}\pill{PATRONES}
\end{center}
\end{frame}}

\end{document}
